% ======================================================================
% Ontology of Continua -- Core 1.3
% Cross-K module: crossk_k5_k6.tex
% Cross-level relations between K5 and K6
% ======================================================================

\subsubsection{Межуровневая структура K5–K6}
\label{sec:crossk-k5-k6}

Переход $K_5 \rightarrow K_6$ обозначает возникновение когнитивного
континуума из нейронного континуума. Если $K_5$ поддерживает электрическую
возбудимость, распространение спайков и сетевую синхронизацию, то $K_6$
вводит:
\begin{itemize}
  \item концептуальные оси и семантические различения,
  \item когнитивные потенциалы (уверенность, релевантность, значимость),
  \item пороги когерентности и противоречивости,
  \item структурированные потоки внимания, памяти и интерпретации,
  \item циклы понимания, прогнозирования и обучения,
  \item внутренний когнитивный временной масштаб $\tau(K_6)$.
\end{itemize}

Этот переход делает возможным появление смысла, вывода, репрезентации и
построения моделей — способностей, непредставимых одной только нейронной
динамикой.

% ----------------------------------------------------------------------
\subsubsection{1. Задействованные уровни и их роли}

\paragraph{K5 (нейронный континуум).}
$K_5$ предоставляет:
\begin{itemize}
  \item распространение и синхронизацию спайков,
  \item динамику каналов и оси мембранного напряжения,
  \item пороги возбудимости $\Theta_{\mathrm{exc}}$,
  \item нейронные потоки $J_5$ и циклы распространения,
  \item сетевую связность и осцилляторные режимы.
\end{itemize}

Однако $K_5$ не включает:
\begin{itemize}
  \item семантическую структуру,
  \item абстрактные репрезентации,
  \item пороги, основанные на когерентности или противоречии,
  \item устойчивые концептуальные аттракторы.
\end{itemize}

\paragraph{K6 (когнитивный континуум).}
$K_6$ вводит:
\begin{itemize}
  \item концептуальные оси $A^{c}$, определяющие различия между смыслами,
  \item потенциалы $P^{c}$, описывающие заметность, силу убеждений,
        релевантность,
  \item пороги $\Theta^{c}$ для когерентности, противоречия,
        когнитивной перегрузки и коллапса,
  \item потоки $J^{c}$: внимание, память, интерпретация, аргументация,
  \item циклы $C^{c}$: цикл внимания, понимания, обучения,
  \item меру $k(K_6)$ для описания когнитивной когерентности.
\end{itemize}

Так $K_6$ становится первым уровнем, способным моделировать свои собственные
внутренние состояния.

% ----------------------------------------------------------------------
\subsubsection{2. Совместные и наследуемые оси и пороги}

\paragraph{Наследуемые оси.}
Нейронные синхронные шаблоны служат подложкой для концептуальных осей:
\[
A_{\mathrm{sync}}(K_5) \rightarrow A^{c}(K_6).
\]
Осцилляторные и связные паттерны становятся семантическими различиями.

\paragraph{Новые оси.}
$K_6$ вводит:
\begin{itemize}
  \item $A^{c}$ — оси концептуальной дифференциации,
  \item $A_{\mathrm{belief}}$ — степень уверенности,
  \item $A_{\mathrm{value}}$ — градиенты релевантности или полезности,
  \item $A_{\mathrm{model}}$ — координаты внутренней модели.
\end{itemize}

\paragraph{Новые пороги.}
$\Theta^{c}$ включает:
\begin{itemize}
  \item \textbf{пороги когерентности} — минимальная совместимость концептов,
  \item \textbf{пороги противоречия} — верхний предел напряжения в модели,
  \item \textbf{пороги перегрузки} — пределы потока внимания или памяти,
  \item \textbf{пороги коллапса} — $\Theta_{\mathrm{collapse}}$, где
        $\Omega(K_6)$ становится пустой при превышении противоречий.
\end{itemize}

Нарушение $\Theta_5$ приводит к коллапсу возбудимости и, следовательно,
делает невозможным $K_6$.

% ----------------------------------------------------------------------
\subsubsection{3. Межуровневые потоки, циклы и напряжения}

\paragraph{Потоки.}
Когнитивные потоки $J^{c}$ расширяют нейронные потоки $J_5$:
\[
J^{c} = (J_{\mathrm{attention}},\ J_{\mathrm{memory}},\
        J_{\mathrm{interpretation}},\ J_{\mathrm{argument}}).
\]
Примеры:
\begin{itemize}
  \item сдвиги внимания, модулирующие нейронную синхронизацию,
  \item консолидация памяти, формирующая нейронные аттракторы,
  \item интерпретационные потоки, переоценочные $P^{c}$,
  \item аргументативные потоки, перестраивающие концептуальные оси.
\end{itemize}

\paragraph{Циклы.}
$K_6$ содержит устойчивые циклы:
\begin{align*}
C_{\mathrm{attention}} &: \text{ориентация} \rightarrow \text{фокус}
    \rightarrow \text{деактивация} \rightarrow \text{переориентация},\\[4pt]
C_{\mathrm{understanding}} &: \text{предсказание} \rightarrow \text{вход}
    \rightarrow \text{исправление} \rightarrow \text{обновлённая модель},\\[4pt]
C_{\mathrm{learning}} &: \text{активация} \rightarrow \text{ошибка}
    \rightarrow \text{коррекция} \rightarrow \text{стабилизация}.
\end{align*}

\paragraph{Напряжение.}
Межуровневое напряжение:
\[
T_{5,6} = f(\Theta^{c} - P^{c},\ J^{c} - J_5,\ A^{c} - A_{\mathrm{sync}}).
\]
Если $T_{5,6}$ превышает размерностный порог, концептуальная когерентность
разрушается и континуум возвращается в не-семантический нейронный режим.

% ----------------------------------------------------------------------
\subsubsection{4. Условия рождения и исчезновения между уровнями}

\paragraph{Рождение \texorpdfstring{$K_6$}{K_6}.}
Когнитивный континуум появляется, когда:
\[
T_5 > \Theta_{5,\mathrm{dim}}
\quad\text{и}\quad
A^{c} \in M_6 \setminus A(K_5),
\]
что соответствует появлению концептуальной дифференциации, не сводимой к
нейронным осям.

Расширение области состояний:
\[
\Omega(K_6) = \Omega(K_5) \cup \Delta\Omega_{\mathrm{cog}}.
\]

\paragraph{Распространение гибели.}
Если $\Theta_5$ нарушена (потеря возбудимости), $K_6$ становится невозможной.

Если $\Theta^{c}$ нарушена (противоречие, перегрузка, коллапс), $K_6$
исчезает, но $K_5$ остаётся корректной:
\[
\Omega(K_6) \rightarrow \varnothing,
\qquad
\Omega(K_5) \text{ не изменяется}.
\]

% ----------------------------------------------------------------------
\subsubsection{5. Концептуальные примеры}

\paragraph{Формирование семантической оси.}
Повторяющаяся синхронизация в $K_5$ стабилизируется в $K_6$ как
концептуальные различения, создавая оси:
\[
A^{c} = \{\text{объект vs фон},\ \text{истина vs ложь},\
         \text{причина vs следствие},\ \ldots\}.
\]

\paragraph{Внутренние модели.}
Прогностические циклы создают устойчивые структуры в $\Omega(K_6)$, представляя
гипотезы, понятия и ожидания.

\paragraph{Когнитивный коллапс.}
Если:
\[
T_{5,6} > \Theta_{\mathrm{collapse}},
\]
пространство концептов теряет когерентность, и $\Omega(K_6)$ схлопывается до нуля.

Следовательно, переход K5→K6 можно в сжатой форме описать как возникновение
концептуальной структуры, смысла и внутренних моделей из нейронной динамики
и синхронизации.
